% Section 9.2: Which operations need a metric? | Part: 1 | Depends on: sec:vectors to sec:duality | Status: done | Original lines: 1723-1752
%!TEX root = ../../main.tex
\subsection{Which operations need a metric?}\label{sec:whichmetric}

The divisions drawn in Sections~\ref{sec:metric}--\ref{sec:duality} can be gathered into one table; each row is justified where the
object is introduced.

\begin{center}
\small
\renewcommand{\arraystretch}{1.3}
\begin{tabular}{@{}p{4.9cm}p{2.2cm}p{7.2cm}@{}}
\toprule
\textbf{Operation / object} & \textbf{Metric?} & \textbf{What it needs instead}\\
\midrule
$\langle\omega,v\rangle$, tensor contraction, $\otimes$ & no & $V$, $V^*$\\
$\wedge$, $\iota_v$, $\bigwedge V$, pull-back $\phi\pb$ & no & $V$, $V^*$\\
$\dd$, $\Lie_X$, integration of forms, Stokes & no & smooth structure (and orientation for $\int$)\\
$\operatorname{div}_\mu X$ (Section~\ref{sec:gauss}) & no & a volume form $\mu$ (a metric supplies one)\\
$\flat,\ \sharp$, $\nabla f=(\dd f)^\sharp$ & yes & $g$\\
Hodge star $\hodge$, dual $\dual A=A\pss^{-1}$ & yes & $g$ and an orientation\\
$a\times b$, curl & yes & $g$, orientation, and $n=3$\\
Geometric product $ab$, $a\lc A$, rotors & yes & $Q$ (the metric)\\
Levi-Civita connection, geodesics, curvature & yes & $g$ (a bare connection needs only the manifold)\\
\bottomrule
\end{tabular}
\end{center}

\begin{remark}[dual basis versus reciprocal frame]\label{rem:reciprocal}
The \emph{dual basis} $e^i\in V^*$ is defined without a metric by $e^i(e_j)=\delta^i{}_j$. The \emph{reciprocal frame} of geometric algebra is the vector $e^i_\sharp=g^{ij}e_j\in V$ with $e^i_\sharp\cdot e_j=\delta^i{}_j$; it is the dual basis transported by $\sharp$. For an orthonormal frame with $e_i^2=\varepsilon_i=\pm1$ it is simply $e^i_\sharp=\varepsilon_ie_i$.
% REWRITTEN: e_i^\sharp -> e^i_\sharp (one notation for the reciprocal frame; docs/NOTATION.md) This is the reciprocal frame
used by the vector derivative $\vd=e^i_\sharp\,\partial_i$ in Section~\ref{sec:ga-calculus}.
\end{remark}
